\cotefile{93}
\scanpage{151}
donnée d'un hom de faisceaux
\[
 P \xrightarrow{\lambda_{\deleted{\unsure{P|E}}}} P^\infty_{S/T}(E)/\delta E \tag*{(26)}
\]
satisfaisant la condition
\[
 \deleted{\lambda(x+\unsure{\xi|\zeta})=\lambda(x)+d^\infty_{S/T}(\xi)\qquad x\in P(\cdot),\ \xi\in E(\cdot).} \tag*{\deleted{(27)}}
\]
\[
 \lambda(x+\xi)=\lambda(x)+p(d^\infty_{S/T}(\xi))\qquad x\in P(\cdot),\ \xi\in E(\cdot). \tag*{(27)}
\]
Il faut encore exprimer que la trivialisation du $P^\infty_{S/T}(E)/E$-torseur obtenu est compatible avec les stratifications. Comme $P$ est \deleted{le} obtenu des sections horizontales de $P^\infty_{S/T}(P)$, on voit que \deleted{\unsure{cette|une}}
% slanted
\margin{plutôt lemme à dégager pour un hom de $P^\infty_{S/T}(P)$ dans un torseur stratifié quelc.}
\deleted{condition \unsure{de|du}} \inserted{cela} signifie qu'on a
\[
 \lambda(P)\subset\bigl(P^\infty_{S/T}(E)/E\bigr)^{\natural} \,, \tag*{(28)}
\]
où le \deleted{\unsure{donné|demi}} signe $\natural$ désigne le sous faisceau des sections horizontales.

\deleted{En} \uline{En résumé}, \deleted{\unsure{une struc|une donnée}} si $P$ est un $E$-torseur, une stratification sur $E$ rel. à $T$, compatible avec la stratification $\delta=\delta^\infty_E$ donnée sur $E$, équivaut à un hom de faisceaux (26), satisfaisant (27) et (28) ; cette dernière condition pouvant d'ailleurs s'exprimer en précisant la donnée (26) par
% Two vertical emphasis strokes alongside the four preceding lines.
\[
 P\xrightarrow{\lambda}\bigl(P^\infty_{S/T}(E)/\delta E\bigr)^\natural . \tag*{(26$'$)}
\]
Cette description est fonctorielle en $E$, et compatible aux changements de base.
% archive sheet number: 150
\scanpage{152}
3) \uline{Relations entre stratifications et connexions intégrables sur torseurs affines.}

Si $P$ est un torseur sous $E$, $E$ muni d'une stratification $\delta^\infty_E=\delta$ rel. à $T$, toute stratification de $P$ \inserted{(compatible à $\delta$)} définit évidemment une connexion (compatible à $d^0_E$).

\deleted{Il est immédiat que l'hom canonique (24), \unsure{au sens de|vu ci-dessus}, est compatible avec cette association.}

Il convient encore de préciser quelle est l'hom $d^0_P:P\longrightarrow E\otimes\Omega_{S/T}$ correspondant à la description (17). Pour l'expliciter, notons que l'inclusion
\[
 P^\infty_{S/T}(E)^+\hookrightarrow P^\infty_{S/T}(E)
\]
fournit, par composition, un isomorphisme additif
\[
 P^\infty_{S/T}(E)^+\simeq P^\infty_{S/T}(E)/\delta E \,, \tag*{(29)}
\]
% slanted
\margin{Définir les sections \deleted{\unsure{hom|hor}} \uline{intégrables} de $P^\infty_{S/T}(E)^+$ comme celles correspondant aux \deleted{\unsure{éléments}} \inserted{sections} de $\bigl(P^\infty_{S/T}(E)/\delta E\bigr)^\natural$, et introduire la notation $P^\infty_{S/T}(E)^+$ [Mais Attention, $P^\infty_{S/T}(E)^+\hookrightarrow P^\infty_{S/T}(E)$ \uline{n'est pas} compatible avec stratifications}
l'hom correspondant

\diagram{\begin{tikzcd}[ampersand replacement=\&,column sep=large]
{(29')}\quad P^\infty_{S/T}(E) \arrow[r,"p",no head] \arrow[rr,bend right=15,no head] \& P^\infty_{S/T}(E)/\delta E \arrow[r,"\sim"] \& P^\infty_{S/T}(E)^+
\end{tikzcd}}

étant donné par
\[
 \deleted{d^+_E}\ \varphi\longmapsto\varphi-\delta.\varepsilon(\varphi)\,. \tag*{\deleted{\unsure{(30)}}}
\]
Donc l'hom $\lambda$ de (26) s'interprète comme un
\[
 \lambda^+:P\longrightarrow P^\infty_{S/T}(E)^+ \,, \tag*{(30)}
\]
\deleted{la condition $\lambda^+$} la condition (27) s'écrivant alors
\[
 \lambda^+(x+\unsure{\xi|\zeta})=\lambda^+(x)+\left[\,\underbrace{d^\infty_{S/T}(\xi)-\delta^\infty_E(\xi)}_{d^+_{S/T}(\xi)}\,\right]. \tag*{(31)}
\]
% archive sheet number: 151
\scanpage{153}
Réduisant à valeurs dans $P^1_{S/T}(E)^+=E\otimes\Omega_{S/T}$, on trouve un hom
\[
 d^0_P:P\longrightarrow E\otimes\Omega_{S/T},\qquad d^0_E(x)=\lambda^+(x)\ \mathrm{mod}\ \ldots \tag*{(32)}
\]
satisfaisant
\[
 d^0_P(x+\unsure{\xi|\zeta})=\unsure{d^0_P|d_P}(x)+d^0_E(\xi)\,,
\]
compte tenu de (23).

% Double vertical emphasis line.
On vérifie que (32) n'est autre que l'hom qui définit la connexion de $P$ associée à sa stratification, \deleted{\unsure{et}} que cette connexion est intégrable, enfin que sa classe dans $H^1_{\mathrm{DR}}(S,E)$ n'est autre que l'image de celle de \deleted{\unsure{$(P,\delta_P)$}} $(P,\lambda)$ dans $H^1_{\mathrm{Strat}}(\unsure{S|A},E_{\mathrm{Strat}})$ par l'homomorphisme canonique (24).
% slanted
\margin{\uline{faire} la \uline{vérification}}

\uline{Corollaire} Supposons que $\unsure{J|T}$ soit de car. nulle et $\Omega_{S/T}$ plat [p. ex. $S$ form\textsuperscript{t} lisse sur $T$] de sorte que (24) est un \uline{isomorphisme} [\enspace]. \deleted{Alors} \inserted{Soit} \deleted{$P$ un \unsure{torseur}} \inserted{$E$ qu. coh. sur $S$} avec stratification $\delta=\delta^\infty_E$, $P$ un torseur sous $E$. Alors il y a correspondance biunivoque entre les stratifications de $P$ compatibles avec celle de $E$, et les connexions intégrables sur $P$ compatibles avec celle de $E$,
\margin{\uline{NB} Sous les hyp. faites, les stratifications sur $E$ correspondent bijectivement aux conn. intégrables sur $E$ \ldots..}
Donc aussi, entre les applications (26) \deleted{\unsure{satisfa}} \inserted{$P\xrightarrow{\lambda}P^\infty_{S/T}(E)/\deleted{\unsure{D}}\delta E$} [ou encore, $\lambda^+:P\longrightarrow (P^\infty_{S/T}(E)^+)^{\deleted{\unsure{\natural}}}$] satisfaisant les conditions (27) (28) [\deleted{\unsure{im|ou}, \unsure{moyen|équivalent}} (31),] \deleted{et \unsure{(28)}, \unsure{à|i.e.}} \inserted{et à valeurs dans $P^\infty_{S/T}(E)^{+\natural}$}
% archive sheet number: 152
\scanpage{154}
\deleted{fait $\lambda=p\lambda^+$}], et les hom (17) $P\xrightarrow{d_P}E\otimes\Omega_{S/T}$ satisfaisant les conditions (18), en associant à $\lambda^+$ sa réduction dans $P^1_{S/T}(E)^+=E\otimes\Omega_{S/T}$.

\uline{NB} Prenant $P$ trivial sur $E$, ceci prouve que l'hom \deleted{de réduction} de réduction $P^\infty_{S/T}(E)^+\longrightarrow P^1_{S/T}(E)^+$ induit un \deleted{\unsure{biject|biject.|bijection}} \uline{isomorphisme de faisceaux}
\[
 P^\infty_{S/T}(E)^{+\natural}\xrightarrow{\sim}{\unsure{Z|\mathbb Z}}^{1}(E\otimes\Omega^\bullet_{S/T})\,,
\]
en prenant au deuxième m\textsuperscript{b} comme stratification sur $P^\infty_{S/T}(E)^+$ celle déduite de la strat. can. de $P^\infty_{S/T}(E)/\delta E$ par transport de structure à l'aide de (29). Ainsi, \uline{une} 1-\uline{forme fermée sur} $S/T$ \inserted{à \uline{coeff. dans} $E$} \uline{définit un jet d'ordre} $\infty$ \inserted{$\lambda^+$} \uline{de} $E$ \uline{sur} $S/T$, satisfaisant la condition (28). A la forme exacte $d^0_E(\unsure{\xi|\zeta})$ est associé ainsi \deleted{\unsure{$\lambda(\xi)=d_E(\xi)$}} \deleted{\unsure{$\lambda^+(\xi)$}}, $\lambda(\xi)=d^\infty_E(\xi)$ i.e. $\lambda^+(\xi)=d^\infty_E(\xi)-\delta^\infty_E(\xi)=d^+_E(\xi)^{\unsure{+|1}}$.

4) \uline{Equations de Picard Fuchs}.

Supposons encore qu'on ait $E$ sur $S$, avec stratification $\delta^\infty_E=\delta$. Supposons donné de plus une suite exacte de modules qu. coh.
\[
 0\longrightarrow E'\mathrel{\overset{u}{\text{---}}}E\mathrel{\overset{v}{\text{---}}}E''\mathrel{\text{---}}0 . \tag*{(\unsure{33|38})}
\]
On ne suppose pas que $E'$ soit un sous-module stratifié. \deleted{\unsure{Supposons}} Supposons donné un
% archive sheet number: 153
\scanpage{155}
% letterhead
% Printed upside down at the lower right of the scanned sheet.
\begin{center}
INSTITUT\\
DES HAUTES ETUDES SCIENTIFIQUES\\
91 - BURES-SUR-YVETTE\\
928 48-53\quad -\quad 928 59-77
\end{center}
% archive sheet number: 154
\scanpage{156}
torseur $P'$ sous $E'$, posons
\[
 P=P'\mathbin{\overset{E'}{\times}}E . \tag*{(34)}
\]
La donnée d'une stratification sur $P$ compatible avec celle de $E$, grâce au n° 2, revient à la donnée d'un hom
\[
 \lambda':P'\longrightarrow \bigl(P^\infty_{S/T}(E)/\delta E\bigr)^\natural \tag*{(35)}
\]
satisfaisant
\[
 \lambda'(x'+{\unsure{\xi|\zeta}}')=\lambda'(x')+p\unsure{\bigl(}d^\infty_E(u\xi')\qquad
 \begin{aligned}x'&\in P'(\cdot)\\ \xi'&\in E'(\cdot).\end{aligned} \tag*{(36)}
\]
On peut aussi interpréter $\lambda'$ comme un
\[
 \lambda'{}^+:P'\longrightarrow\bigl(P^\infty_{S/T}(E)^+\bigr)^\natural \,, \tag*{(35$'$)}
\]
\margin{(cf. $\overline{\mathrm{II}}$ structure).}
satisfaisant
\[
 \lambda'{}^+(x'+\xi')=\lambda'{}^+(x')+
 \underbrace{d^+_E(u\xi')}_{=\,d^\infty_E(u\xi')-\delta^\infty_E(u\xi')}\,. \tag*{(36)$'$}
\]
\deleted{et une condition d'\uline{intégrabilité}, signifiant}

\deleted{Notons alors} Notons alors la suite exacte \inserted{de p-faisceaux}
\[
 \unsure{0|\cdots}\mathrel{\text{---}}P^\infty_{S/T}(E')\mathrel{\overset{u}{\text{---}}}P^\infty_{\unsure{S/T|S+}}(E)\xrightarrow{\deleted{\unsure{p}}v}P^\infty_{S/T}(E'')\longrightarrow 0 \,, \tag*{(37)}
\]
\inserted{compatibles aux stratifications}, d'où, grâce à \deleted{\unsure{(35)}} \inserted{(35) et (36)}, une section canonique
\[
 \varphi=\varphi(\unsure{\lambda|P|\lambda'})\in\Gamma(S,P^\infty_{S/T}(E'')/\unsure{v}\delta(E))\simeq
 \Gamma(S,P^\infty_{S/T}(E'')^+/v\delta u(E')) \,, \tag*{(38)}
\]
\deleted{\unsure{i.e. un|En outre}} \inserted{En fait}, \deleted{\unsure{q|$\varphi$}} on a \deleted{\unsure{un}} même
\[
 \varphi\in\Gamma\bigl(S,(P^\infty_{S/T}(E'')/\delta\unsure{(|v}E)^\natural\bigr)\simeq
 \Gamma\bigl(S,P^\infty_{S/T}(E'')^{+\natural}/\unsure{\delta u(E')}\bigr). \tag*{(38$'$)}
\]
% Right ends of (38), (38') meet the edge of the sheet.
% archive sheet number: 155
\scanpage{157}
On l'appelle le \deleted{\unsure{secti|section}} \uline{jet de Picard-Fuchs défini par} $P'$ \deleted{\unsure{$\delta_P$}} \uline{et la stratification $\sigma$ donnée} \deleted{\unsure{de}} $P=P'\mathbin{\overset{E'}{\times}}E$. Ce n'est pas tt à fait un jet \inserted{intégrable} de $E''$, mais son \unsure{déguisement} \unsure{affine|affaibli} : c'est \unsure{là|le|à-dire} un tel jet, défini modulo \deleted{\unsure{une sec}} \inserted{l'image \inserted{par $vd^\infty_Eu$} d'une} section de \deleted{\unsure{$P^\infty_{S/T}(E')$}} $E'$.

Soit $F$ un module \inserted{qu. coh.} sur \deleted{$E$} $S$, et
\[
 D:E''\longrightarrow F \tag*{(39)}
\]
un op différentiel \deleted{\unsure{d'or}} rel. / $T$, définissant donc un hom \inserted{linéaire} continu : \deleted{\unsure{hom}}
\[
 u_D:P^\infty_{S/T}(E'')\longrightarrow F . \tag*{(40)}
\]
On dit que \deleted{$u_D$} $D$ est un \uline{opérateur de Picard-Fuchs}, \uline{relativement à la suite exacte} (33), si l'application linéaire correspondante (40) s'annule sur $v\delta(E)$, i.e. si $u_D$ se factorise en un homomorphisme
\[
 \overline u_D:P^\infty_{S/T}(E'')/v\delta(E)\longrightarrow F . \tag*{(41)}
\]
\deleted{Soit alors} Si alors $(P',\sigma)$ sont comme dessus, d'où un $\varphi$ (38), on peut former $\overline u_D(\varphi)\in\deleted{\unsure{H}}\Gamma(S,F)$ ; \unsure{on posera|en posant|on pose}
\[
 \deleted{\unsure{u_D(\varphi)=}}\ \mu_D(P,\sigma)=\overline u_D(\varphi(P,\sigma))\in\Gamma(S,F)\,, \tag*{(42)}
\]
% archive sheet number: 156
\scanpage{158}
on l'appellera la \uline{section de Manin de} $F$ associée à l'op. de Picard-Fuchs $D$ et aux données $(P',\sigma)$.

Équation de Pic. Fuchs universelle : à valeur dans le pro-faisceau $P^\infty_{S/T}(E'')/\unsure{u|v}\delta u(E')$ ; elle redonne la section $\varphi$ du début.
% Two pairs of horizontal rules below this note.

Lorsque $E$ est seulem\textsuperscript{t} muni d'une connexion intégrable \deleted{\unsure{, de}} \inserted{rel. à $T$} (au lieu d'une stratification), on voit que les connexions intégrables sur $P=P'\mathbin{\overset{E'}{\times}}E$ correspondent aux hom
\[
 P'\xrightarrow{d_{P'}}\deleted{\unsure{E'}}E\otimes\Omega_{S/T} \tag*{(43)}
\]
satisfaisant
\[
 \left\{\begin{aligned}
 d_{P'}(x'+{\unsure{\xi|\zeta}}')&=d_{P'}(x')+d^0_E(u(\xi'))\\
 d_{P'}(P')&\subset Z^1(E\otimes\Omega_{S/T}),\ \text{i.e.}\quad d^{\unsure{1|2}}_E\circ d_{P'}=0 .
 \end{aligned}\right. \tag*{(44)}
\]
Dans ce cas, on a une variante de la notion d'équation \inserted{de} Picard-Fuchs \inserted{pour (30)}, en considérant les opérateurs différentiels
\[
 E\otimes\Omega^1_{S/T}\xrightarrow{\Delta\deleted{\unsure{D}}}F \tag*{(45)}
\]
nuls sur \deleted{l'image de $E'\otimes\Omega_{S/T}$} \inserted{$d^0_E(u(E'))$}, \deleted{i.e. les opérateurs différentiels sur} [qu'il suffira de prendre modulo ceux de la forme $\Delta'\circ d^1_E$, \deleted{$E''\otimes\Omega_{S/T}\longrightarrow$} où \deleted{$\Delta'$} $E\otimes\Omega^2_{S/T}\mathrel{\overset{\Delta'}{\text{---}}}F$ un autre op. diff]. Comme $d_{P'}$ (43) fournit \inserted{loc\textsuperscript{t}} une section de $E\otimes\Omega_{S/T}$ modulo une section
% archive sheet number: 157
\scanpage{159}
de la forme $d^0_E(u({\unsure{\xi|\zeta}}'))$, $\xi'\in E'(\cdot)$, on voit que $\Delta$ définit une section
\[
 \unsure{\nu|v}_\Delta(P',d_{P'})\in\Gamma(S,F). \tag*{(46)}
\]
Ici la situation universelle est obtenue en prenant pour $F$ le préfaisceau
\[
 F=P^\infty_{S/T}(E\otimes\Omega^{1}_{S/T})/\operatorname{Im}P^\infty_{S/T}(E')\,,
\]
où $P^\infty_{S/T}(E')\longrightarrow P^\infty_{S/T}(E\otimes\Omega_{S/T})$ provient de l'op. diff. composé (de \uline{degré} 1) \inserted{$d^0_Eu$:} $E'\mathrel{\overset{u}{\text{---}}}E\mathrel{\overset{d^0_E}{\text{---}}}E\otimes\Omega_{S/T}$.

Revenons au cas où $E$ admet une stratification, \deleted{et $P$ muni d'une stratification, compatible avec} et \unsure{prenant|prenons} $d^0_E$ déduit d'icelle ; \inserted{soit $(P',\deleted{\unsure{hom}}\sigma)$} \deleted{pour ramener} (comme on \unsure{disait|didait} \unsure{au|que} \unsure{n°|où}. \deleted{les sections de la forme} On ramène les sections $\nu_\Delta$ aux $\mu_D$ de la façon suivante. À tout \deleted{\unsure{$\nu_\Delta$}} \inserted{$\Delta$} est associé un opérateur de Picard-Fuchs, \deleted{\unsure{savoir pour la}} pour (30), en prenant le composé
\[
 D_0:E\xrightarrow{d^0_E}E\otimes\Omega_{S/T}\xrightarrow{\Delta}F \,, \tag*{(47)}
\]
et en passant au quotient, \deleted{\unsure{ce qui}} \inserted{en $D:E''\longrightarrow F$,} ce qui est possible puisque par hyp. sur $\Delta$\quad $\Delta\circ d^0_E$ est nulle sur $E'$. Ceci posé, on trouve que
\[
 \nu_\Delta=\mu_D \,, \tag*{(48)}
\]
% slanted
\margin{faire la vérification}
ce qui donne la réduction annoncée.
% archive sheet number: 158
\scanpage{160}
Dans le cas favorable envisagé à la fin du n° 3, ($\unsure{J|T}$ de car. nulle, $\Omega_{S/T}$ plat), on peut de $=$ exprimer les $\mu_D$ en termes des $\unsure{\nu|v}_\Delta$. Pour ceci, il suffit de voir que, pour $S$ affine, chaque $D$ provient bien d'un $\Delta$. Ceci résulte du fait suivant (lui-$=$ cas particulier d'un fait bien plus \unsure{générale|général} sur les op. diff) : \deleted{Tout op.}

\uline{Lemme}. Sous les conditions précédentes ($S$ car. nulle, $\Omega_{S/T}$ plat), \deleted{tout} \inserted{les} op diff $E\longrightarrow F$ qui sont tels que $P^\infty_{S/T}(E)\longrightarrow F$ s'annule sur $\delta(E)$, sont exactement ceux \deleted{qui \unsure{se}} de la forme $\Delta\circ d^0_E$, avec $\Delta:E\otimes\Omega_{S/T}\mathrel{\text{---}}F$ un op différentiel. L'indétermination est \unsure{exactement|exactem\textsuperscript{t}} \deleted{\unsure{l'op}} par un op diff. \deleted{\unsure{$D', E\otimes\Omega_{S/T}^2$|$\Delta'\circ d^1_E:E\otimes\Omega_{S/T}$}} \inserted{$\Delta'\circ d^1_E$}, où $\Delta':E\otimes\Omega_{S/T}\longrightarrow F$.

\deleted{Ceci résulte du fait que \unsure{$E\longrightarrow\Omega E$|$E\longrightarrow\Omega$}.}
% archive sheet number: 159
\scanpage{161}
Modules stratifiés et \deleted{Connexions intégrables et} équations de Picard-F
% The right-hand cut in the cover meets the last letter.
% archive sheet number: 160
\scanpage{162}
\mspage{\unsure{1|I}}
1) \uline{Filtration du complexe de DeRham d'un schéma relatif \inserted{lisse}, et connexion de Gauss-Manin} \deleted{\unsure{en dim.}} (d'après Katz-Oda).

\diagram{\begin{tikzcd}[ampersand replacement=\&,row sep=large]
X \arrow[d,"{\substack{f\ \text{lisse [ou seulement diff. lisse,}\\\text{pour rendre exacte la}\\\text{suite (1)]}}}"]\\
S \arrow[d,"g"]\\
T
\end{tikzcd}}

On a donc une suite exacte
\[
 \deleted{\text{On a donc}}\ 0\longrightarrow f^*(\Omega^1_{S/T})\longrightarrow\Omega^1_{X/T}\longrightarrow\Omega^1_{X/S}\longrightarrow0 . \tag*{(1)}
\]
d'où une filtration $\mathrm{Fil}_{X/S}$ de l'objet gradué $\Omega^\bullet_{X/T}=\bigwedge^{*}\Omega^1_{X/T}$, avec gradué associé
\[
 \mathrm{Gr}^p_{X/S}(\Omega^\bullet_{\deleted{\unsure{S|X}}X/T})\simeq f^*(\Omega^p_{S/T})\otimes\Omega^\bullet_{X/S}[-p] \tag*{(2)}
 % Faint pencil brackets above the first factor on the left and beneath the shifted factor on the right.
\]
On constate que la filtration respecte les opérateurs \unsure{différentiel|différentiels} de $\Omega^\bullet_{X/T}$, et qu'il en est de $=$ de l'isomorphisme (2), \inserted{compte tenu que, grâce à la conn. intégrable de $f^*(\Omega^p_{S/T})$ rel : $X/S$, $f^*(\Omega^p_{S/T})\otimes\Omega^\bullet_{X/S}$ \unsure{se trouve|étant|est} bien muni d'une différentielle.} [D'ailleurs, comme $f$ est plat, on a
\[
 f^*(\Omega^p_{S/T})\simeq Lf^*(\Omega^p_{S/T})
\]
et comme $\Omega^\bullet_{X/S}[-p]$ est plat, le $\otimes$ de (2) peut se remplacer par $\overset{\mathbb L}{\otimes}$, donc on a
\[
 \mathrm{Gr}^p_{X/S}(\Omega^\bullet_{X/T})\simeq\mathbb L f^*(\Omega^p_{S/T})\overset{\mathbb L}{\otimes}\Omega^\bullet_{X/S}[-p] \,, \tag*{(2$'$)}
\]
dans la mesure où cela veut dire qqe chose].

La filtration sur $\Omega^\bullet_{X/T}$ définit $\mathbb R f_*(\Omega^\bullet_{X/T})$ comme objet d'une catégorie dérivée filtrée, (filtration encore notée $\mathrm{Fil}^\bullet_{X/S}$, d'où $\mathrm{Gr}^\bullet_{X/S}$). On
% archive sheet number: 161
\scanpage{163}
aura alors
\[
 \mathrm{Gr}^p_{X/S}\mathbb R f_*(\Omega^\bullet_{X/T})\simeq\mathbb R f_*(\mathrm{Gr}^p_{X/S}(\Omega^\bullet_{X/T})) .
\]
Tenant compte de (2$'$), et supposant $f$ qu. cpt. qu. sép., il résulte \inserted{alors} de la formule de projection [dans un cadre pas très catholique, - il faut faire une p\textsuperscript{t}ite vérification]
\[
 \mathrm{Gr}^p_{X/S}\mathbb R f_*(\Omega^\bullet_{X/T})\simeq\Omega^p_{S/T}\overset{\mathbb L}{\otimes}\mathbb R f_*(\Omega^\bullet_{X/S})[-p] \tag*{(3)}
\]
On a \deleted{en particulier des opérateurs} une différentielle
\[
 d_1^{(p)}:\mathrm{Gr}^p_{X/S}\mathbb R f_*(\Omega^\bullet_{X/T})\mathrel{\text{---}}\mathrm{Gr}^{p-1}_{X/S}\mathbb R f_*(\Omega^\bullet_{X/T})[1] \,, \tag*{(4)}
\]
qui donne, compte tenu de l'expression explicite (3) :
\[
 d_1^{(p)}:\Omega^p_{S/T}\overset{\mathbb L}{\otimes}\mathbb R f_*(\Omega^\bullet_{X/S})[-p]\mathrel{\text{---}}\Omega^{p\unsure{+|-}1}_{S/T}\overset{\mathbb L}{\otimes}\mathbb R f_*(\Omega^\bullet_{X/S})[-p] \,, \tag*{(5)}
\]
d'où, en translatant par $p$ :
\[
 \deleted{\unsure{\beta|p}}\quad d_1^{(p)}[p]:\Omega^p_{S/T}\overset{\mathbb L}{\otimes}\mathbb R f_*(\Omega^\bullet_{X/S})\mathrel{\text{---}}\Omega^{p+1}_{S/T}\overset{\mathbb L}{\otimes}\mathbb R f_*(\Omega^\bullet_{X/S}). \tag*{(6)}
\]
On a bien sur
\[
 d_1^{(p+1)}d_1^{(p)}=0 . \tag*{(7)}
\]
Faisant $p=0$ dans (6), on trouve
\[
 d_1^{(0)}:\mathbb R f_*(\Omega^\bullet_{X/S})\mathrel{\text{---}}\Omega^1_{S/T}\otimes\mathbb R f_*(\Omega^\bullet_{X/S}) \,, \tag*{$(\unsure{6_0'|6_0})$}
\]
i.e. une connexion sur $\mathbb R f_*(\Omega^\bullet_{X/S})$. \deleted{La formule (7) prouve que c'est une} On l'appellera la \uline{connexion de Gauss-Manin}. (Il semble bien que ce soit la $=$ que celle que j'ai définie dans [\enspace] !). La formule (7) montre qu'\uline{elle est intégrable}.
% slanted
\margin{! Attention, il faut prouver une compatibilité \unsure{des|de} $d_1$ avec diff. de $\Omega^\bullet_{S/T}$, pour \unsure{exploiter|expliciter} \unsure{comme on doit|comment on déd.|comme dit} (7) !}
% archive sheet number: 162
\scanpage{164}
\deleted{Lorsque les} Bien sur, les différentielles \inserted{(6)} définissent
\[
 d_1^{p,q}[p]:\mathscr H^q\bigl(\Omega^p_{S/T}\otimes\mathbb R f_*(\Omega^\bullet_{X/S})\bigr)\longrightarrow
 \mathscr H^q\bigl(\Omega^{p+1}_{S/T}\otimes\mathbb R f_*(\Omega^\bullet_{X/S})\bigr) \,, \tag*{(\deleted{\unsure{6}}8)}
\]
qui, pour $\Omega^p_{S/T}$ plat sur $S$, p. ex $S$ formellement lisse sur $T$, \deleted{\unsure{def}} s'écrit aussi
\[
 d_1^{p,q}[p]:\Omega^p_{S/T}\otimes\mathbb R^q f_*(\Omega^\bullet_{X/S})\longrightarrow
 \Omega^{p+1}_{S/T}\otimes\mathbb R^q f_*(\Omega^\bullet_{X/S}) \,, \tag*{(8$'$)}
\]
qui pour $p=0$ définit une connexion intégrable sur les $\mathbb R^q f_*(\Omega^\bullet_{X/S})$, dite conn. de \uline{Gauss-Manin}.

Posons pour un moment, pour simplifier,
\[
 K^\bullet=\mathbb R f_*(\Omega^\bullet_{X/T}) \,, \tag*{(9)}
\]
qu'on considère comme complexe filtré comme on l'a dit. On aura donc
\[
 \mathbb R\Gamma_S(K^{\unsure{\bullet|'}})\simeq\mathbb R\Gamma_X(\Omega^\bullet_{X/T})\simeq\mathbb RH_{\mathrm{DR}}(X/T). \tag*{(10)}
\]
La filtration de $K^\bullet$ permet, d'après un \unsure{sorte|sorite|sortie} général, d'écrire une suite spectrale
\[
 \mathbb R\Gamma_S(K^\bullet)\Longleftarrow E^{\unsure{p,q|pq}}_2=\mathbb H^p(S,\alpha\longmapsto\mathscr H^{\alpha+q}(\mathrm{Gr}^\alpha(K^\bullet))) \,; \tag*{(11)}
\]
les différentielles définissant les complexes $\alpha\mathrel{\unsure{\longmapsto|\mapstochar\relbar\mapstochar}}\ldots$ étant la différentielle évidente $d_1$ provenant de la filtration. Tenant compte de (3) \inserted{et (10)}, on trouve donc la suite spectrale.
% archive sheet number: 163
